% !TEX TS-program = lualatex

% ------------------------------------------------- %
% -- DO NOT MODIFY. FILE GENERATED AUTOMATICALLY -- %
% ------------------------------------------------- %

\documentclass{tutodoc}

\usepackage{../preamble.cfg}


\begin{document}

% -- AUTOMATICALLY GENERATED UGLY CODE - START -- %

\subsection{LaTeX}
\label{products-latex}

\subsubsection{Basic use}

To use a color from a palette, use \tdoccodein{text}+\palUse{<name>}{<index>}+ where \tdoccodein{text}+<name>+ is the standard palette name (without prefix), and \tdoccodein{text}+<index>+ is the color number (ranging from 1 to 10). For example, \tdoccodein{text}+\palUse{GistHeat}{8}+ is the eighth color of the \tdoccodein{text}+GistHeat+ palette, an {\xcolor} format color that can be easily used as shown in the following compilable example.

\begin{tdoccode}{latex}
\documentclass{article}

\usepackage{palettes}
\usepackage{tikz}

\begin{document}

\textcolor{\palUse{GistHeat}{8}}{\bfseries Colored text.}

Representation of the color palette.

\begin{tikzpicture}
  \foreach \i in {1,...,10} {
    \fill[\palUse{GistHeat}{\i}]
      (1.25*\i - 1, 0) rectangle (1.25*\i, 1);
  }
\end{tikzpicture}

\end{document}
\end{tdoccode}

\subsubsection{Creating palettes from scratch}

For creating new palettes manually, the following high-level commands are available.

\begin{enumerate}
\item
\tdoccodein{latex}+\palCreateFromRGB+ creates a palette by entering it as an array-like variable, while \tdoccodein{latex}+\palCreateFromNames+ works with named colors.

\item
\tdoccodein{latex}+\palSize{<name>}+ returns the palette size (useful for loops, for example).

\end{enumerate}

The following example demonstrates the \tdoccodein{latex}+\palCreateFromRGB+ and \tdoccodein{latex}+\palCreateFromNames+ commands (we don't have put \tdoccodein{latex}+\usepackage{palettes}+).

\begin{tdoccode}{latex}
\usepackage[svgnames]{xcolor}

\palCreateFromRGB{MyRGBPal}{
  {0.0, 0.0, 0.0},
  {0.4, 0.0, 0.2},
  {0.8, 0.2, 0.0},
}

\palCreateFromNames{MyNameUsePal}{
  YellowGreen,
  green!60!black,
  LimeGreen!80,
}
\end{tdoccode}

\begin{tdocnote}
All built-in palettes are created using the \tdoccodein{latex}+\palCreateFromRGB+ macro.
\end{tdocnote}

A lower-level approach is also available through the following commands.

\begin{enumerate}
\item
\tdoccodein{latex}+\palNew{<name>}+ initializes a new (empty) palette.

\item
\tdoccodein{latex}+\palAddNames{<name>}{<color-using-names>}+ appends a color defined with named colors to the palette.

\item
\tdoccodein{latex}+\palAddRGB{<name>}{<r>, <g>, <b>}+ appends an \tdoccodein{text}+RGB+ color to the palette, where \tdoccodein{text}+<r>+, \tdoccodein{text}+<g>+, and \tdoccodein{text}+<b>+ are decimal values ranging from 0 to 1.

\end{enumerate}

The following example demonstrates the flexibility offered by these low-level commands.

\begin{tdoccode}{latex}
\usepackage[svgnames]{xcolor}

\palNew{LowLevelPal}
\palAddNames{LowLevelPal}{IndianRed}
\palAddRGB{LowLevelPal}{0.0, 0.0, 0.0}
\palAddNames{LowLevelPal}{green!60!black}
\palAddRGB{LowLevelPal}{0.8, 0.2, 0.0}
\end{tdoccode}

\subsubsection{Creating palettes from existing ones}

Building new palettes by transforming existing ones can be achieved using the \tdoccodein{latex}+\palCreateFromPal+ command, which has the signature \tdoccodein{latex}+\palCreateFromPal{<new-name>}[<options>]{<existing-name>}+. The following example shows how to do this (all options are used).

\begin{tdoccode}{latex}
\palCreateFromPal{BlackbodyTransformed}[
  extract = {1, 3, 6, 9},
  shift   = 1,
  reverse
]{Blackbody}
\end{tdoccode}

\begin{tdoctip}
\tdoccodein{latex}+\palCreateFromPal{<new-name>}{<existing-name>}+ build a copy of an existing palette.
\end{tdoctip}

\subsubsection{Retrieving the internal definition of a palette}

The internally stored definition of a palette named \tdoccodein{text}+MyPal+, for example, is \tdoccodein{latex}+\g_palette_MyPal_seq+ which is a {\LaTeX3} variable (keep in mind the pattern \tdoccodein{latex}+\g_palette_PaletteName_seq+).

\begin{tdocnote}
Variables of type \tdoccodein{latex}+\g_palette_PaletteName_seq+ are not used internally to retrieve the colors themselves; they are only there for technical reasons related to the development process of new palettes via {\LaTeX}.
\end{tdocnote}

% -- AUTOMATICALLY GENERATED UGLY CODE - END -- %

\end{document}
